% Physical unit 016; accepted Session 02.
\GordanUnitStart{016}{printed page 11}
\GordanTopFolio{11}
In quam spectant rationes hae:
\[
 \delta_1\int_{l_1}^{l_2}dl\sqrt{1-e^2\cos^2l+\cos^2l\,{d\psi\over dl}\cdot{d\psi\over dl}}=0
\]
\[
\begin{aligned}
&\delta_2\int_{l_1}^{l_2}dl\sqrt{1-e^2\cos^2l+\cos^2l\,{d\psi\over dl}\cdot{d\psi\over dl}}\\
&=\int_{l_1}^{l_2}\left(
{dl\,\delta\!\bigl({d\psi\over dl}\bigr)\delta\!\bigl({d\psi\over dl}\bigr)\cos^2l[1-e^2\cos^2l]
\over \{1-e^2\cos^2l+\cos^2l\,{d\psi\over dl}\cdot{d\psi\over dl}\}^{1\frac12}}
\right)>0.
\end{aligned}
\]
Quae linea, quoniam propinquis lineis brevior est, nominetur brevissima.

Atque, ne forte error inducatur hoc nomine, eam, qua omnino nulla brevior cogitari potest, omnium brevissimam dicamus, quae pars brevissimae intelligenda est, dum haec ipsa est pars geodeticae.

\begin{center}\S\ 2.\end{center}

Cum integralis:
\[
 \int_{l_1}^{l_2}\left(
{dl\,\delta\!\bigl({d\psi\over dl}\bigr)\delta\!\bigl({d\psi\over dl}\bigr)\cos^2l[1-e^2\cos^2l]
\over \{1-e^2\cos^2l+\cos^2l\,{d\psi\over dl}\cdot{d\psi\over dl}\}^{1\frac12}}
\right)
\]
quodque elementum positivum esse constet, totum integrale positivum sit, necesse est. Cui consentaneum esse videatur, omnem lineam geodeticam vel in tota extensione esse brevissimam.

Quod non ita est, quia Cap. II \S\ 7 demonstravimus, nostrae aequationis vim tantum hactenus stare, quoad circulus \gq{$-l_0$} tangatur linea geodetica.

\begin{center}\S\ 3.\end{center}

Quod igitur ex aequatione integrali, quam supra proposuimus, cognosci non potest, quatenus linea geodetica, tacto etiam circulo \gq{$-l_0$}, brevissima sit; regulae universales statuantur oportet. Quod consequemur inducentes, loco variabilis integrationis $l$, aliam variabilem $\psi$ integrationis, quam pro ratione recedentis ab exordio puncti $\Pi$ perpetuo crescere, manifestum est.

Proinde quaerendum est, quae sit maxima longitudo $\psi$, usque ad quam perpetuo
\[
 \delta_2\int_{\psi_1}^{\psi_2}d\psi\sqrt{(1-e^2\cos^2l)l'^2+\cos^2l}
\]
